% QUESTAO
Podemos calcular o determinante de uma matriz $A$, de ordem $3$, por meio de um dispositivo denominado {\it regra de Sarrus}. Nessa regra, à direita da matriz, repetimos suas duas primeiras colunas
\[\det A = \begin{array}{|ccc|cc}
a_{11} &  a_{12}  &  a_{13}  &  a_{11}  &  a_{12}\\
a_{21} &  a_{22}  &  a_{23}  &  a_{21}  &  a_{22}\\
a_{31} &  a_{32}  &  a_{33}  &  a_{31}  &  a_{32}
\end{array} \]
Em seguida, realizamos multiplicações conservando os sinais dos produtos obtidos no sentido da diagonal principal e mudando os sinais dos produtos obtidos no sentido da diagonal secundária. Assim,
se $A = \left[ \begin{array}{ccc}
    -1  &   4  &   3\\
    0   &  2  &  -3\\
    1  &  -2  &   5
\end{array} \right]$, então $\det A$ é:

% ALTERNATIVAS
$\det A = -22$
$\det A = 45$
$\det A = -25$
$\det A = 30$
$\det A = -18$


